\chapter{Podsumowanie i wnioski}
\label{ch:podsumowanie}

\section{Wnioski}
\label{sec:wnioski}

Celem pracy było zaprojektowanie, zaimplementowanie i~wdrożenie na~klastrze rozproszonego
programu symulacyjnego rozwiązującego trójwymiarowe równanie przewodnictwa ciepła w~języku Chapel oraz przeprowadzenie
systematycznej analizy jego poprawności i~skalowalności. Wszystkie postawione cele szczegółowe
zostały zrealizowane, a~zebrane wyniki empiryczne potwierdzają tezę pracy: dla rozważanej klasy
obliczeń stencilowych o~czynniku ograniczającym wydajność decyduje nie moc obliczeniowa, lecz
przepustowość pamięci na~poziomie węzła oraz koszt komunikacji między węzłami~\cite{roofline}.

Najważniejsze wnioski można ująć następująco.

\paragraph{Poprawność i~równoważność wariantów.} Wykazano, że rozwiązanie różnicowe daje wynik niezależny
od~liczby lokacji, na~które rozdzielono dziedzinę, pod~warunkiem prawidłowej wymiany warstw
brzegowych. Niezmiennik sumy pola w~stanie końcowym okazał się prostym, lecz skutecznym
kryterium kontrolnym. Potwierdzono również, że dwa warianty pętli czasowej, \emph{ping-pong}
oraz wariant z~zamianą wskaźników, są numerycznie identyczne i~wydajnościowo równoważne,
ponieważ zamiana tablic odbywa się w~czasie stałym i~nie wiąże się z~kopiowaniem danych.
Jednocześnie wykazano, że pominięcie wymiany halo, choć znacząco skraca czas wykonania, prowadzi
do~wyniku \emph{cicho błędnego}; obserwacja ta podkreśla, że w~programowaniu rozproszonym sama
szybkość nie jest miarą poprawności.

\paragraph{Skalowanie wątkowe.} Zrównoleglenie wątkowe w~obrębie węzła nasyca się bardzo szybko.
Dla kostki $1000^3$ optimum przypada na~cztery wątki, a~maksymalne przyspieszenie nie przekracza
około $1{,}55\times$ mimo szesnastu dostępnych wątków; wydajność równoległa spada wówczas
poniżej $0{,}1$. Zachowanie to jest charakterystyczne dla obliczeń ograniczonych przepustowością
pamięci: po~wysyceniu pasma dokładanie kolejnych wątków nie zwiększa wydajności, a~jedynie
podnosi koszt synchronizacji i~rywalizacji o~zasoby.

\paragraph{Skalowanie po~węzłach.} Przyspieszenie względem liczby węzłów zatrzymuje się na~poziomie
około $2{,}4\times$ i~pozostaje daleko od~zależności liniowej, a~jego przebieg jest niemonotoniczny: najlepszy czas uzyskano przy~sześciu węzłach. Źródłem niemonotoniczności jest nierównomierny
podział kostki $1000^3$ dla ,,niewygodnych'' dzielników liczby lokacji, spotęgowany wzrostem
wolumenu wymiany halo, który rośnie w~przybliżeniu liniowo z~liczbą węzłów. W~połączeniu
z~ograniczoną przepustowością łącza \SI{1}{\giga\bit\per\second} czyni to obliczenia
ograniczonymi komunikacją.

\paragraph{Skalowanie rozmiaru zadania.} Wraz ze~wzrostem krawędzi kostki rośnie przepustowość
(z~$0{,}117$ do~$2{,}722$ miliarda aktualizacji komórek na~sekundę), a~udział komunikacji maleje
z~$0{,}96$ do~$0{,}45$. Wynika to bezpośrednio z~faktu, że objętość zadania rośnie z~sześcianem
krawędzi, podczas gdy powierzchnia wymiany halo, jedynie z~jej kwadratem. Większe zadania
lepiej wykorzystują klaster; małe są zdominowane przez komunikację i~marnują dostępne zasoby.

\paragraph{Wybór warstwy transportu.} Praca pokazała praktyczne znaczenie doboru conduitu GASNet.
Wariant AMUDP, choć prosty i~pozbawiony zależności, ulega awarii typu \texttt{ECONGESTION}
pod~obciążeniem typu \emph{incast} i~przy stratach pakietów, ponieważ nie zapewnia kontroli
przepływu~\cite{incast}. Wariant mpi-conduit oparty na~bibliotece MPICH i~protokole TCP przetrwa te same
warunki dzięki kontroli przepływu i~retransmisji, kosztem nieco niższej wydajności, lecz
z~zachowaniem poprawności i~stabilności. Dla docelowego środowiska wybór ten okazał się
decydujący.

\paragraph{Produktywność a~wydajność w~języku Chapel.} Wysoki poziom abstrakcji języka Chapel, globalny widok tablic, gotowy rozkład \texttt{StencilDist} wraz z~automatyczną wymianą warstw
brzegowych oraz zwięzłe iteratory równoległe, pozwolił zapisać poprawny kod rozproszony
w~sposób nieporównanie prostszy niż przy jawnym użyciu MPI~\cite{chapel-overview,chapel-book}. Abstrakcja ta nie zniosła jednak
fundamentalnego ograniczenia komunikacyjnego: koszt wymiany halo pozostaje widoczny w~wynikach
i~decyduje o~granicy skalowania. Wniosek ten jest zgodny z~intuicją, żaden model programowania
nie zmienia fizyki sieci, a~zarazem stanowi ważną wskazówkę praktyczną: produktywność
i~wydajność są w~tym przypadku zagadnieniami komplementarnymi, nie zaś zamiennymi.

\section{Ograniczenia}
\label{sec:ograniczenia}

Uzyskane wyniki należy interpretować w~świetle kilku ograniczeń, które jednocześnie wyznaczają
naturalne kierunki dalszych prac.

Po~pierwsze, środowiskiem docelowym był niewielki, jednorodny klaster połączony siecią
Ethernet o~przepustowości \SI{1}{\giga\bit\per\second}. Łącze to jest wąskim gardłem
i~w~znacznej mierze determinuje obserwowane granice skalowania po~węzłach. Na~infrastrukturze
z~szybką siecią o~małych opóźnieniach, na~przykład InfiniBand, należałoby oczekiwać istotnie lepszego
skalowania oraz przesunięcia punktu, w~którym program przestaje być ograniczony komunikacją.
Weryfikacja tej hipotezy wymaga jednak dostępu do~odpowiedniego sprzętu.

Po~drugie, praktyczną granicą jest maksymalny rozmiar zadania: dziedzin istotnie większych
niż~$1000^3$ nie udało się zmieścić w~zasobach pamięciowych i~czasowych klastra.

Po~trzecie, świadomie ograniczono się do~jawnej metody różnic skończonych na~siatce regularnej
z~arytmetyką podwójnej precyzji. Metody niejawne, które są bezwarunkowo stabilne, lecz wymagają
rozwiązywania układów równań, siatki niestrukturalne, techniki nakładania komunikacji na~obliczenia
oraz akceleracja na~procesorach graficznych pozostają poza zakresem pracy. Każde z~tych
rozszerzeń stanowi obiecujący kierunek: w~szczególności nakładanie wymiany halo na~obliczenia
we~wnętrzu poddziedziny mogłoby częściowo ukryć koszt komunikacji, a~przejście na~schemat
niejawny zmieniłoby charakterystykę obliczeniową problemu.

Po~czwarte, badania koncentrowały się na~skalowaniu silnym, przy ustalonym rozmiarze zadania i~rosnącej
liczbie jednostek równoległości. Analiza skalowania słabego, w~którym rozmiar zadania rośnie
proporcjonalnie do~liczby lokacji, dostarczyłaby komplementarnego obrazu i~jest naturalnym
uzupełnieniem przedstawionych wyników.

Mimo tych ograniczeń praca dostarcza spójnego, empirycznie ugruntowanego obrazu zachowania
rozproszonego programu stencilowego w~języku Chapel oraz jednoznacznie identyfikuje komunikację
i~przepustowość pamięci jako czynniki decydujące o~jego skalowalności.

\subsection{Kierunki dalszych prac}
\label{sec:dalsze-prace}

Zidentyfikowane ograniczenia wyznaczają naturalną mapę dalszych badań. Najbardziej obiecującym
kierunkiem jest \textbf{nakładanie komunikacji na~obliczenia}: wymianę warstw brzegowych można
rozpocząć asynchronicznie i~w~tym czasie liczyć wnętrze poddziedziny, które nie zależy od~danych
sąsiadów, ujawniając komunikację dopiero przy aktualizacji komórek przylegających do~granicy.
Technika ta potrafi w~znacznym stopniu ukryć koszt wymiany halo, a~w~języku Chapel jej wyrażenie, za~pomocą zadań asynchronicznych i~jawnego rozdzielenia wnętrza od~brzegu poddziedziny, jest naturalne.

Drugim kierunkiem jest \textbf{zmiana warstwy sprzętowej}: powtórzenie pomiarów na~sieci
o~niskim opóźnieniu i~wysokiej przepustowości, na~przykład InfiniBand, oraz z~conduitem GASNet dobranym
do~takiej sieci pozwoliłoby ustalić, w~jakim stopniu obserwowane granice skalowania wynikają
z~samego algorytmu, a~w~jakim z~ograniczeń konkretnego łącza \SI{1}{\giga\bit\per\second}.
Komplementarnym eksperymentem byłoby \textbf{skalowanie słabe}, w~którym rozmiar zadania rośnie
proporcjonalnie do~liczby lokacji; pozwoliłoby ono oddzielić koszty stałe od~kosztów rosnących
wraz z~rozmiarem problemu.

Trzecim kierunkiem są \textbf{rozszerzenia numeryczne i~sprzętowe}: przejście na~schematy niejawne,
które są bezwarunkowo stabilne, lecz wymagają rozwiązywania układów równań i~globalnej komunikacji innego
typu, obsługa siatek niestrukturalnych oraz akceleracja jądra obliczeniowego na~procesorach
graficznych, którą Chapel wspiera poprzez odpowiednie lokacje i~atrybuty pętli równoległych. Każde
z~tych rozszerzeń zmieniłoby charakterystykę obliczeniową problemu i~pozwoliłoby zweryfikować, na~ile
wnioski niniejszej pracy zachowują ważność poza wąską klasą jawnych schematów stencilowych.

Wreszcie, wartościowym uzupełnieniem o~charakterze inżynierskim byłoby \textbf{pełne
zautomatyzowanie potoku pomiarowego}, od~budowy i~dystrybucji, przez zbieranie samoopisujących
się dzienników, aż po~generowanie wykresów, tak aby cały cykl badawczy był odtwarzalny jednym
poleceniem. Wprowadzone w~tej pracy elementy, czyli samoopisujące się logi oraz skrypty dystrybucji
i~analizy, stanowią solidny fundament takiego potoku.
